\chapter{相关理论与技术基础}
\echapter{Theoretical Background and Technical Foundations}

承接第一章对研究问题、总体文献格局和研究空缺的概括，本章围绕本文方法设计所依赖的关键概念和技术脉络展开。具体而言，本章进一步回答“这些范式和机制为什么与本文方法有关”，为后续通信框架设计、机制修正和实验诊断提供理论基础。

\section{多智能体强化学习问题建模}
\esection{Problem Formulation of Multi-Agent Reinforcement Learning}

多智能体系统（Multi-Agent Systems, MAS）的研究最早可追溯至 20 世纪 70 年代的分布式人工智能（Distributed Artificial Intelligence, DAI）。随着深度学习与强化学习的结合，多智能体强化学习（Multi-Agent Reinforcement Learning, MARL）逐渐成为解决复杂协作决策问题的重要范式。与单智能体强化学习相比，MARL 不仅需要解决状态估计与策略优化问题，还必须处理多智能体之间因目标耦合、观测受限与策略共适应带来的额外复杂性。

协作型 MARL 问题通常可建模为分布式部分可观测马尔可夫决策过程（Dec-POMDP）。设系统中共有 $N$ 个智能体，状态空间为 $\mathcal{S}$，智能体 $i$ 的局部观测空间与动作空间分别为 $\mathcal{O}_i$ 与 $\mathcal{A}_i$，联合动作记为 $\mathbf{a}=(a_1,\ldots,a_N)$，共享奖励函数记为 $r(s,\mathbf{a})$。智能体的目标是在长期交互过程中学习联合策略 $\pi(\mathbf{a}|\mathbf{o})$，使折扣累计回报最大。

与单智能体场景相比，Dec-POMDP 中的核心难点主要体现在两个方面。其一，单个智能体的观测通常不足以恢复全局状态，因此仅凭局部信息往往难以形成最优联合决策；其二，其他智能体策略在训练过程中同步变化，会使任一智能体所感知到的环境呈现强非平稳性，从而进一步加大信用分配、价值估计和策略收敛的难度。本文关注的显式通信问题，本质上正是在这一背景下提出：当局部观测不足、但协作又依赖队友即时意图时，如何通过受控的信息交互打破“信息孤岛”，成为实现高水平协同智能的关键。

\section{协作决策基础：从独立学习到 CTDE}
\esection{Collaborative Decision Making: From Independent Learning to CTDE}

在多智能体系统的早期研究中，一个最直接的想法是让每个智能体独立使用强化学习算法进行训练。Tan 的经典工作指出，独立学习（Independent Q-Learning, IQL）虽然实现简单，但会将队友策略变化视为环境变化的一部分，从而导致严重的非平稳性问题，最终使得学习过程难以稳定收敛\cite{tan1993iql}。换言之，IQL 把本应显式处理的“协作关系”隐含进了环境噪声中，因此在合作任务上的表现通常明显弱于真正的协作训练方案。

另一端的极端方案是完全中心化的训练与执行，即训练和执行都假设全局状态可得。该方案在理论上能够直接优化联合策略，但其执行阶段依赖全局信息，因而并不适用于大多数真实的部分可观测场景。为平衡全局协同与局部可执行性，集中式训练、分布式执行（Centralized Training and Decentralized Execution, CTDE）逐渐成为协作型 MARL 的主流范式。

在 CTDE 框架下，训练阶段允许利用额外的集中信息，例如全局状态、联合动作或其他智能体中间表示；而执行阶段则要求每个智能体仅依据自身局部观测和历史隐藏状态进行决策。值分解方法是这一范式下的重要分支。VDN 将团队总价值函数分解为各智能体局部价值之和\cite{sunehag2018vdn}，QMIX 则进一步通过单调混合网络刻画联合价值与局部价值之间的非线性关系，在保持分布式执行的同时提升了协作表达能力\cite{rashid2018qmix}。在此基础上，QTRAN 尝试突破单调性约束带来的表达限制\cite{son2019qtran}，QPLEX 则进一步扩展了 IGM 框架下的联合价值表达能力\cite{wang2021qplex}。这类方法在训练阶段能够实现隐式协同，但其执行阶段仍只能依赖局部观测，因此在高度部分可观测环境中往往容易出现策略迟滞或信息利用不足的问题。
在异构多智能体任务中，国内研究进一步将价值函数分解与通信学习机制结合，说明信息交互也可以作为增强值分解协作能力的重要补充\cite{du2024heterocomm_cn}。

除值分解方法外，基于策略梯度的 CTDE 框架也得到了广泛发展。COMA 通过集中式价值评估器（centralized critic）与反事实基线（counterfactual baseline）缓解多智能体信用分配问题\cite{foerster2018coma}；MADDPG 则将行动者-评论家（actor-critic）框架推广到混合协作与竞争环境\cite{lowe2017maddpg}。Proximal Policy Optimization Algorithms 提出的 PPO 截断更新机制为稳定策略优化提供了重要基础\cite{schulman2017ppo}；The Surprising Effectiveness of PPO in Cooperative Multi-Agent Games 进一步表明，在适当的集中式价值评估与训练细节配合下，多智能体近端策略优化（Multi-Agent Proximal Policy Optimization, MAPPO）能够在协作任务中取得极强表现\cite{yu2021mappo}。沿着这一方向，Trust Region Policy Optimisation in Multi-Agent Reinforcement Learning 进一步系统讨论了 HATRPO/HAPPO 一类 trust-region 风格多智能体策略优化方法\cite{kuba2022happo}，FACMAC 则展示了中心化策略梯度与分解式 Critic 结合的潜力\cite{peng2021facmac}，MAT 进一步说明了序列建模范式也能够成为强有力的协作决策主干\cite{wen2022mat}。这一点也影响了后续通信研究的实现范式：越来越多的方法不再从较弱主干直接学习通信，而是建立在相对稳定的行动者-评论家主干或值分解主干之上，再讨论通信结构如何提供额外协同收益。例如，Context-aware Communication for Multi-agent Reinforcement Learning 同时适配行动者-评论家框架与值函数分解框架（value-based framework）\cite{li2024cacom}；Asynchronous Cooperative Multi-Agent Reinforcement Learning with Limited Communication 也建立在更明确的协作决策主干之上\cite{dolan2025asyncomarl}。本文选择 MAPPO 作为主干策略框架，正是因为它在保持 CTDE 范式优势的同时，兼具较强的稳定性与较好的结构兼容性，适合作为后续通信机制研究的强基线。

\section{信息交互拓扑的演进}
\esection{The Evolution of Communication Topology}

当环境部分可观测程度较高时，仅依靠 CTDE 框架中训练阶段的集中信息并不能彻底弥补执行阶段的信息缺失。显式通信机制的核心目标，正是在执行期允许智能体交换消息，以补充局部观测所无法覆盖的状态信息或协作意图。

早期通信方法通常采用全连接或广播式结构。例如，CommNet 通过在多个智能体之间传播连续消息向量，实现了端到端可微的联合训练\cite{sukhbaatar2016commnet}；DIAL 则进一步讨论了可微通信信道在协作任务中的训练方式\cite{foerster2016dial}。这类方法验证了“可学习通信”本身的可行性，但也暴露出一个明显问题：当智能体规模扩大时，全连接或广播式通信会导致信息冗余、计算复杂度上升以及消息语义被平均化。

为此，研究者开始探索如何通过选择性交互来降低通信负载并提高消息质量。注意力机制的引入，使得智能体能够像人类一样根据当前任务上下文动态决定“听谁的”和“对谁说话”。TarMAC 利用目标化通信思想，通过注意力选择与签名机制实现了更具方向性的消息传递\cite{das2019tarmac}；ATOC 则通过注意力单元决定智能体是否应与邻近智能体形成临时通信组，从而实现动态变化的局部交互拓扑\cite{jiang2018attentional}。这些方法的共同特点是：不再默认所有智能体都应对所有邻居平均通信，而是把交互关系本身建模为一个状态依赖的选择问题。
这一方向在近年进一步演化为更明确的结构化通信学习。Learning Structured Communication 通过层级化拓扑组织通信分组与组间路由，强调通信边并非固定存在，而可以随任务需求动态形成\cite{sheng2023lsc}。国内综述工作也把基于通信的 MARL 作为独立脉络进行整理，并指出部分可观测环境下通信机制与协同控制之间具有直接联系\cite{wang2022commmarl_cn}。

从图结构视角看，通信拓扑本质上也是一种可学习结构，消息路由、边选择与结构搜索正在成为通信研究的重要方向。进一步地，MAGI 从图信息瓶颈角度讨论了如何在保留有效协同信息的同时压缩消息冗余并增强鲁棒性\cite{ding2024magi}，IDEAL 则说明在图消息传递过程中显式区分智能体身份，有助于提升通信表达能力并抑制行为同质化\cite{du2024ideal}。从任务图建模角度出发的协作子任务行为发现研究也表明，识别智能体间动态交互集合可以帮助提升局部与全局协作效率\cite{li2024csp_cn}。这也说明，通信研究的重点已经从“能不能通信”逐步转向“应该以怎样的拓扑通信”。

\section{通信内容与信息内涵的深化}
\esection{Deepening the Semantics of Communication}

通信机制的有效性不仅取决于“和谁通信”，还取决于“传递什么信息”。如果消息本身缺乏可解释语义或与任务目标关联过弱，那么即使路由结构设计得当，通信也可能退化为高熵噪声。为解决这一问题，研究者逐步从抽象特征交换转向更具任务意义的队友建模与意图推断。

MAIC 是这一方向上的代表性工作之一。它通过去中心化的队友建模与互信息约束，让每个智能体能够推测队友的决策偏好，并生成具有激励作用的定制消息，从而显式调整其他智能体的价值评估\cite{yuan2022maic}。与仅传播隐藏状态或连续向量相比，这种方法强调消息应当服务于协作目标本身，即通信不只是传递信息，更是在影响队友的决策倾向。
近年的工作也开始从消息表征本身提高通信质量。Learning Multi-Agent Communication with Contrastive Learning 利用对比学习增强消息表示的一致性与区分性，说明通信学习不一定只依赖更复杂的路由，也可以从表征质量入手改善下游协同\cite{lo2024learning}。Language Grounded Multi-agent Reinforcement Learning with Human-interpretable Communication 则进一步尝试把智能体通信空间对齐到人类可解释语言表示，使消息既服务任务，又具备更强的语义可读性\cite{li2024languagegroundedmarl}。
与此相近，面向智能体的语义通信研究强调通信目标不应停留在符号传输本身，而应关注语义是否能够以期望方式影响接收端行为\cite{zhang2022semanticcomm_cn}。

这类方法的共同趋势是：通信内容不再被视为无语义的中间特征，而是被提升为对队友意图、局部计划甚至奖励感知的显式建模。对于部分可观测协作任务而言，这种“心智模型化”的通信设计往往更有机会在复杂环境中带来稳定协同。从更细的机制角度看，Unveiling Decision Intention for Cooperative Multi-Agent Reinforcement Learning 进一步把“决策意图”显式作为分析对象，强调理解和生成队友意图本身也是协作透明性的重要组成部分\cite{zhang2025decisionintention}。强化学习可解释性研究同样提示，若要将学习型策略用于复杂协作场景，需要进一步解释环境、任务与策略之间的关系\cite{liu2023xrl_cn}。

然而，这类方法也通常具有更高的建模复杂度与训练成本。互信息目标、队友模型误差以及潜变量表示稳定性，都会为优化带来新的难点。因此，在追求更丰富语义的同时，如何控制通信模块的训练风险，仍然是一个未被彻底解决的问题。

\section{交互时机与通信成本控制}
\esection{Timing Control and Communication Cost}

现实场景中的通信往往受到带宽、延迟和能耗约束，因此“何时通信”本身也是一个关键问题。若智能体在每个时刻都进行高带宽交互，不仅会带来额外计算开销，还可能把不必要的噪声持续注入决策过程。

为实现按需交互，研究者引入了门控机制（gating mechanism）与稀疏通信策略。其核心思想是：在简单状态下保持独立决策，在冲突状态或复杂决策点才消耗通信资源。SchedNet 则把有限带宽下“谁先说、谁来听”的调度问题显式纳入学习过程\cite{kim2019schednet}。这样的设计不仅有助于降低通信成本，也有助于抑制背景噪声。
在此基础上，Model-based Sparse Communication 进一步讨论了如何通过消息估计与调度机制减少不必要的消息发送，从而在保持决策质量的同时压缩信道开销\cite{han2023mbc}。Context-aware Communication for Multi-agent Reinforcement Learning 则强调消息不应只是统一广播，而应根据接收方上下文生成个性化内容，以便在带宽受限场景下提高单位通信的有效性\cite{li2024cacom}。

除通信频率外，通信问题并不只是静态的信息交换问题，还涉及消息是否在正确的时间窗口内被生成、传输和利用。Revisiting Communication Efficiency in Multi-Agent Reinforcement Learning from the Dimensional Analysis Perspective 进一步指出，通信效率不仅与是否通信有关，也与接收端最终保留了多少真正有用的消息维度有关\cite{sun2025drmac}。
进一步地，Asynchronous Cooperative Multi-Agent Reinforcement Learning with Limited Communication 表明，在异步且有限通信条件下，通信拓扑与消息时机需要被共同建模，这也使“有限成本下的有效通信”成为更贴近真实系统的重要议题\cite{dolan2025asyncomarl}。从协作增益形成机理看，Situation-Dependent Causal Influence-Based Cooperative Multi-Agent Reinforcement Learning 则进一步强调，不同状态下智能体之间的影响强度并不一致，只有在真正存在互相影响的关键状态中，通信或协同信号才更有可能转化为有效决策收益\cite{du2024scic}。

从本文的研究角度看，这些工作提供了两个重要启发。其一，通信不能简单理解为“加一个消息通路”，而应被视为具有成本约束的条件性决策过程；其二，门控是否打开、在什么状态下打开，以及打开后对动作层到底产生了多大影响，都是需要被单独诊断和验证的机制变量。

\section{架构自动化搜索与图结构学习}
\esection{Automated Architecture Search and Graph Structure Learning}

另一条路线试图自动学习通信拓扑，即不再预设固定的通信边和交互模式，而是将通信结构本身作为可优化对象，与策略参数共同训练。在这一思路下，通信图不只是消息传递的载体，也被看作影响协作效率的重要结构变量。相关方法通常通过图建模、边选择或层次化组织，使智能体能够根据任务状态动态形成更适合当前协作需求的交互关系。

这一方向的代表性特点，在于它把“谁与谁通信”从人工设定转变为可学习对象。从研究趋势看，这类方法体现了多智能体通信从“设计一个通信模块”走向“联合学习通信架构”的进一步发展。

不过，架构自动化搜索与图结构学习通常也意味着更大的结构搜索空间、更复杂的优化耦合关系以及更高的训练成本。对于强调稳定性、低风险接入和动作层可解释修正的本文而言，这一路线虽具有重要启发意义，但并不完全符合本文“轻量通信”的设计取向。因此，本文并未继续沿着复杂拓扑自动搜索展开，而是选择在固定而受控的通信框架内，重点研究通信如何稳定、有效地介入动作决策。

\section{与本文工作的关系}
\esection{Relation to This Thesis}

综合上述研究可以看到，现有多智能体通信方法大致沿着四条主线演进：其一，从独立学习走向 CTDE，以缓解非平稳性与信用分配问题；其二，从广播式通信走向注意力与图结构路由，以提升通信拓扑的选择性；其三，从抽象特征交换走向意图建模与互信息约束，以增强消息语义；其四，从始终通信走向门控、稀疏化与延迟建模，以降低通信成本并提升时机匹配性。

尽管这些方向都取得了重要进展，但它们仍留下一个在本文实验中尤为突出的空缺：大量工作关注的是消息能否生成、路由是否合理、门控是否稀疏，却较少直接分析通信是否真正稳定地进入了最终动作决策过程，以及它是否在真正需要干预的状态中改变了动作选择。也就是说，通信模块可能在表示层、路由层甚至门控层都看似“工作正常”，但在动作层依然几乎没有决策影响。

本文的研究正是在这一空缺上展开。与侧重提升通信表达能力、扩大消息带宽或自动搜索复杂拓扑的路线不同，本文更关注轻量、受控且可解释的通信修正：先构建强主干策略，再让通信以残差形式进入有限动作子空间，并通过一系列机制诊断区分通信失败究竟发生在沉默逃逸、路由退化、门控关闭，还是“通信到动作层的作用力”不足。进一步地，本文并不满足于“通信存在”这一表层结论，而是把“队友目标与本地目标冲突”“有效动作扰动”和“融合后决策是否改变”等指标纳入分析，从而把研究问题推进到更细的层面：通信是否在该介入的状态中，沿着正确方向改变最终动作。
这一关注点也与近年关于通信效率与信息利用方式的讨论形成互补。相比于仅从消息发送或接收表示本身讨论通信是否高效，本文进一步把分析推进到动作层介入与决策影响诊断层面，更直接地回答“通信是否真的改变了决策”这一问题。

这一研究视角决定了本文在方法设计上的基本原则：通信内容应尽可能具备明确物理语义；通信路由和门控应受到预算与结构约束；通信更适合作为局部决策的残差修正，而非替代整个策略主干；而真正困难的问题，不只是“如何让通信存在”，而是“如何让通信在保持稳定的同时具有可验证的动作层决策影响价值”。这也正是本文后续算法设计与实验分析的核心出发点。

\section{本章小结}
\esection{Summary}

本章围绕多智能体强化学习中的协作决策与显式通信问题，对相关理论与技术基础进行了梳理。首先，本文从 Dec-POMDP 建模与 CTDE 范式出发，说明了在部分可观测条件下引入集中训练与分布式执行主干的必要性。其次，围绕通信拓扑、通信内容、交互时机与通信成本控制等方向，总结了现有研究从“能否通信”逐步走向“如何高效通信”的演进脉络。最后，结合现有方法在动作层实际介入效果上的不足，明确了本文后续工作的关注重点，即在低通信代价约束下，设计能够稳定作用于动作决策的轻量通信机制，并对其决策影响进行可验证分析。
